%-------------------------
% Resume in Latex
% Author : Jake Gutierrez
% Based off of: https://github.com/sb2nov/resume
% License : MIT
%------------------------
%
% Original source: https://github.com/jakegut/resume
% The layout and macros below are Jake Gutierrez's original, unchanged. Only the
% sample content was replaced with placeholders, by CVEdge
% (https://www.thecvedge.com/cv-format/jakes-resume). Not affiliated with or
% endorsed by the author.
%
% Compile with pdfLaTeX (the default on Overleaf).
%
%------------------------
% MIT License
%
% Copyright (c) 2020 Jake Gutierrez
%
% Permission is hereby granted, free of charge, to any person obtaining a copy
% of this software and associated documentation files (the "Software"), to deal
% in the Software without restriction, including without limitation the rights
% to use, copy, modify, merge, publish, distribute, sublicense, and/or sell
% copies of the Software, and to permit persons to whom the Software is
% furnished to do so, subject to the following conditions:
%
% The above copyright notice and this permission notice shall be included in all
% copies or substantial portions of the Software.
%
% THE SOFTWARE IS PROVIDED "AS IS", WITHOUT WARRANTY OF ANY KIND, EXPRESS OR
% IMPLIED, INCLUDING BUT NOT LIMITED TO THE WARRANTIES OF MERCHANTABILITY,
% FITNESS FOR A PARTICULAR PURPOSE AND NONINFRINGEMENT. IN NO EVENT SHALL THE
% AUTHORS OR COPYRIGHT HOLDERS BE LIABLE FOR ANY CLAIM, DAMAGES OR OTHER
% LIABILITY, WHETHER IN AN ACTION OF CONTRACT, TORT OR OTHERWISE, ARISING FROM,
% OUT OF OR IN CONNECTION WITH THE SOFTWARE OR THE USE OR OTHER DEALINGS IN THE
% SOFTWARE.
%------------------------

\documentclass[letterpaper,11pt]{article}

\usepackage{latexsym}
\usepackage[empty]{fullpage}
\usepackage{titlesec}
\usepackage{marvosym}
\usepackage[usenames,dvipsnames]{color}
\usepackage{verbatim}
\usepackage{enumitem}
\usepackage[hidelinks]{hyperref}
\usepackage{fancyhdr}
\usepackage[english]{babel}
\usepackage{tabularx}
\input{glyphtounicode}


%----------FONT OPTIONS----------
% sans-serif
% \usepackage[sfdefault]{FiraSans}
% \usepackage[sfdefault]{roboto}
% \usepackage[sfdefault]{noto-sans}
% \usepackage[default]{sourcesanspro}

% serif
% \usepackage{CormorantGaramond}
% \usepackage{charter}


\pagestyle{fancy}
\fancyhf{} % clear all header and footer fields
\fancyfoot{}
\renewcommand{\headrulewidth}{0pt}
\renewcommand{\footrulewidth}{0pt}

% Adjust margins
\addtolength{\oddsidemargin}{-0.5in}
\addtolength{\evensidemargin}{-0.5in}
\addtolength{\textwidth}{1in}
\addtolength{\topmargin}{-.5in}
\addtolength{\textheight}{1.0in}

\urlstyle{same}

\raggedbottom
\raggedright
\setlength{\tabcolsep}{0in}

% Sections formatting
\titleformat{\section}{
  \vspace{-4pt}\scshape\raggedright\large
}{}{0em}{}[\color{black}\titlerule \vspace{-5pt}]

% Ensure that generate pdf is machine readable/ATS parsable
\pdfgentounicode=1

%-------------------------
% Custom commands
\newcommand{\resumeItem}[1]{
  \item\small{
    {#1 \vspace{-2pt}}
  }
}

\newcommand{\resumeSubheading}[4]{
  \vspace{-2pt}\item
    \begin{tabular*}{0.97\textwidth}[t]{l@{\extracolsep{\fill}}r}
      \textbf{#1} & #2 \\
      \textit{\small#3} & \textit{\small #4} \\
    \end{tabular*}\vspace{-7pt}
}

\newcommand{\resumeSubSubheading}[2]{
    \item
    \begin{tabular*}{0.97\textwidth}{l@{\extracolsep{\fill}}r}
      \textit{\small#1} & \textit{\small #2} \\
    \end{tabular*}\vspace{-7pt}
}

\newcommand{\resumeProjectHeading}[2]{
    \item
    \begin{tabular*}{0.97\textwidth}{l@{\extracolsep{\fill}}r}
      \small#1 & #2 \\
    \end{tabular*}\vspace{-7pt}
}

\newcommand{\resumeSubItem}[1]{\resumeItem{#1}\vspace{-4pt}}

\renewcommand\labelitemii{$\vcenter{\hbox{\tiny$\bullet$}}$}

\newcommand{\resumeSubHeadingListStart}{\begin{itemize}[leftmargin=0.15in, label={}]}
\newcommand{\resumeSubHeadingListEnd}{\end{itemize}}
\newcommand{\resumeItemListStart}{\begin{itemize}}
\newcommand{\resumeItemListEnd}{\end{itemize}\vspace{-5pt}}

%-------------------------------------------
%%%%%%  RESUME STARTS HERE  %%%%%%%%%%%%%%%%%%%%%%%%%%%%


\begin{document}

%----------HEADING----------
\begin{center}
    \textbf{\Huge \scshape Your Name} \\ \vspace{1pt}
    \small 555-010-0000 $|$ \href{mailto:you@example.com}{\underline{you@example.com}} $|$
    \href{https://linkedin.com/in/yourname}{\underline{linkedin.com/in/yourname}} $|$
    \href{https://github.com/yourname}{\underline{github.com/yourname}}
\end{center}


%-----------EDUCATION-----------
% Education sits first for students and new graduates. Move this section below
% Experience once you have a year or two of full-time work.
\section{Education}
  \resumeSubHeadingListStart
    \resumeSubheading
      {University Name}{City, ST}
      {Bachelor of Science in Computer Science}{Aug. 2022 -- May 2026}
    \resumeSubheading
      {Second School (optional)}{City, ST}
      {Degree or Program}{Aug. 2020 -- May 2022}
  \resumeSubHeadingListEnd


%-----------EXPERIENCE-----------
% Each bullet: what you built or changed, how, and what it did. Add a number
% wherever you have one (users, requests per second, minutes saved, percent).
\section{Experience}
  \resumeSubHeadingListStart

    \resumeSubheading
      {Software Engineering Intern}{May 2025 -- Aug. 2025}
      {Company Name}{City, ST}
      \resumeItemListStart
        \resumeItem{Built [feature] with [language or framework], used by [number] [users or teams]}
        \resumeItem{Cut [metric, e.g. page load time] from [before] to [after] by [what you changed]}
        \resumeItem{Wrote [tests, docs or tooling] that [result, e.g. caught regressions before release]}
      \resumeItemListEnd

% -----------Multiple Positions Heading-----------
%    \resumeSubSubheading
%     {Software Engineer I}{Oct 2014 - Sep 2016}
%     \resumeItemListStart
%        \resumeItem{Apache Beam}
%          {Apache Beam is a unified model for defining both batch and streaming data-parallel processing pipelines}
%     \resumeItemListEnd
%    \resumeSubHeadingListEnd
%-------------------------------------------

    \resumeSubheading
      {Job Title}{Sep. 2023 -- May 2025}
      {Organization Name}{City, ST}
      \resumeItemListStart
        \resumeItem{Action verb, what you did, and the outcome with a number}
        \resumeItem{Keep each bullet to one or two lines}
      \resumeItemListEnd

  \resumeSubHeadingListEnd


%-----------PROJECTS-----------
% For new graduates this section carries the application. Say what the project
% does and which part you built, not just the stack.
\section{Projects}
    \resumeSubHeadingListStart
      \resumeProjectHeading
          {\textbf{Project Name} $|$ \emph{Language, Framework, Database}}{Jan. 2025 -- Present}
          \resumeItemListStart
            \resumeItem{One line on what the project does and who it is for}
            \resumeItem{What you built: [component], using [technology], handling [scale or number]}
          \resumeItemListEnd
      \resumeProjectHeading
          {\textbf{Another Project} $|$ \emph{Language, Framework, Tool}}{Sep. 2024 -- Dec. 2024}
          \resumeItemListStart
            \resumeItem{One line on what the project does and who it is for}
            \resumeItem{A result you can point to: [downloads, users, stars or a measured improvement]}
          \resumeItemListEnd
    \resumeSubHeadingListEnd


%-----------PROGRAMMING SKILLS-----------
% Mirror the wording in the job posting. List what you could be interviewed on.
\section{Technical Skills}
 \begin{itemize}[leftmargin=0.15in, label={}]
    \small{\item{
     \textbf{Languages}{: Language, Language, Language} \\
     \textbf{Frameworks}{: Framework, Framework, Framework} \\
     \textbf{Developer Tools}{: Tool, Tool, Tool} \\
     \textbf{Libraries}{: Library, Library, Library}
    }}
 \end{itemize}


%-------------------------------------------
\end{document}
